Insects exhibit a remarkable repertoire of behaviors that superficially resemble markers of self-awareness in vertebrates:
metacognitive uncertainty monitoring~\cite{Giurfa2013}, caste-appropriate role identity~\cite{SciAdv2024Caste},
tool-use anticipation requiring forward models~\cite{NaturePuzzle2024}, and individual variation in general cognitive
ability (GCA)~\cite{AnimalCog2024}. Yet these phenomena have been studied in isolation, with no unifying computational
account of how a miniature brain could implement the self-referential computations they apparently require.

The central complex (CX), a conserved mid-brain structure across insects, is an increasingly compelling candidate for
a shared substrate. The CX encodes body-centric spatial representations~\cite{Stone2017CX}, regulates circadian
timing, and gates action selection. Critically, it maintains a real-time model of the organism's current state---a
property essential to self-referential behavior. We propose that a predictive coding (PC) architecture
instantiated in the CX, governed by a single \emph{precision} parameter, can account for all four self-awareness
markers simultaneously.

\subsection*{Defining functional self-awareness}

A crucial distinction must be established at the outset. We use the term ``functional self-awareness'' to denote a
computational property: the capacity to maintain a precision-weighted generative model of one's own state, where
self-model prediction errors drive behavior. This is distinct from \emph{phenomenal consciousness}---the presence of
subjective experience or ``what it is like to be''~\cite{Metzinger2003}. We make no claim that bees are conscious;
our model is agnostic to that question. The behaviors we model (opt-out on hard trials, caste-appropriate task
selection, anticipatory tool orientation) can be produced by a purely functional self-model without any accompanying
inner experience~\cite{Barron2016}.

This operationalization follows the Free Energy Principle (FEP)~\cite{Friston2010}: under FEP, agents
minimize expected free energy by maintaining generative models of their environment and themselves. The precision
parameter in FEP quantifies the confidence placed on prediction errors, determining whether sensory surprises
update internal models or are suppressed. Our model inherits this formal apparatus while focusing on the behavioral
signatures precision variation predicts across domains.

\subsection*{The precision-metacognition paradox}

A central and counterintuitive prediction emerges directly from our model: \emph{more intelligent bees should be
worse at metacognitive opt-out tasks}. This ``overconfidence paradox'' arises because high precision, while
improving task performance, simultaneously reduces the uncertainty signal that drives opt-out behavior. We derive
this result analytically and confirm it computationally (metacognition--GCA correlation $r = -0.658$ in simulations
of 500 agents). The prediction is directly falsifiable by correlating GCA battery scores~\cite{AnimalCog2024} with
opt-out accuracy in the same individuals.

\subsection*{Paper overview}

We first formalize the CX predictive coding model and its four behavioral domain functions
(Section~\ref{sec:model}). We then derive the precision--metacognition trade-off analytically and identify the
optimal precision level $\pstar = 0.563$ (Section~\ref{sec:analytical}). Full-scale simulations ($N = 500$ agents
per condition) validate the model's quantitative predictions across four experimental conditions
(Section~\ref{sec:simulations}). We discuss three directly testable empirical predictions and situate the model
within the broader self-awareness and FEP literature (Section~\ref{sec:discussion}).
